% Options for packages loaded elsewhere
\PassOptionsToPackage{unicode}{hyperref}
\PassOptionsToPackage{hyphens}{url}
\documentclass[
  11pt,
]{article}
\usepackage{xcolor}
\usepackage[margin=2.5cm]{geometry}
\usepackage{amsmath,amssymb}
\setcounter{secnumdepth}{-\maxdimen} % remove section numbering
\usepackage{iftex}
\ifPDFTeX
  \usepackage[T1]{fontenc}
  \usepackage[utf8]{inputenc}
  \usepackage{textcomp} % provide euro and other symbols
\else % if luatex or xetex
  \usepackage{unicode-math} % this also loads fontspec
  \defaultfontfeatures{Scale=MatchLowercase}
  \defaultfontfeatures[\rmfamily]{Ligatures=TeX,Scale=1}
\fi
\usepackage{lmodern}
\ifPDFTeX\else
  % xetex/luatex font selection
\fi
% Use upquote if available, for straight quotes in verbatim environments
\IfFileExists{upquote.sty}{\usepackage{upquote}}{}
\IfFileExists{microtype.sty}{% use microtype if available
  \usepackage[]{microtype}
  \UseMicrotypeSet[protrusion]{basicmath} % disable protrusion for tt fonts
}{}
\usepackage{setspace}
\makeatletter
\@ifundefined{KOMAClassName}{% if non-KOMA class
  \IfFileExists{parskip.sty}{%
    \usepackage{parskip}
  }{% else
    \setlength{\parindent}{0pt}
    \setlength{\parskip}{6pt plus 2pt minus 1pt}}
}{% if KOMA class
  \KOMAoptions{parskip=half}}
\makeatother
\usepackage{longtable,booktabs,array}
\newcounter{none} % for unnumbered tables
\usepackage{calc} % for calculating minipage widths
% Correct order of tables after \paragraph or \subparagraph
\usepackage{etoolbox}
\makeatletter
\patchcmd\longtable{\par}{\if@noskipsec\mbox{}\fi\par}{}{}
\makeatother
% Allow footnotes in longtable head/foot
\IfFileExists{footnotehyper.sty}{\usepackage{footnotehyper}}{\usepackage{footnote}}
\makesavenoteenv{longtable}
\usepackage{graphicx}
\makeatletter
\newsavebox\pandoc@box
\newcommand*\pandocbounded[1]{% scales image to fit in text height/width
  \sbox\pandoc@box{#1}%
  \Gscale@div\@tempa{\textheight}{\dimexpr\ht\pandoc@box+\dp\pandoc@box\relax}%
  \Gscale@div\@tempb{\linewidth}{\wd\pandoc@box}%
  \ifdim\@tempb\p@<\@tempa\p@\let\@tempa\@tempb\fi% select the smaller of both
  \ifdim\@tempa\p@<\p@\scalebox{\@tempa}{\usebox\pandoc@box}%
  \else\usebox{\pandoc@box}%
  \fi%
}
% Set default figure placement to htbp
\def\fps@figure{htbp}
\makeatother
% definitions for citeproc citations
\NewDocumentCommand\citeproctext{}{}
\NewDocumentCommand\citeproc{mm}{%
  \begingroup\def\citeproctext{#2}\cite{#1}\endgroup}
\makeatletter
 % allow citations to break across lines
 \let\@cite@ofmt\@firstofone
 % avoid brackets around text for \cite:
 \def\@biblabel#1{}
 \def\@cite#1#2{{#1\if@tempswa , #2\fi}}
\makeatother
\newlength{\cslhangindent}
\setlength{\cslhangindent}{1.5em}
\newlength{\csllabelwidth}
\setlength{\csllabelwidth}{3em}
\newenvironment{CSLReferences}[2] % #1 hanging-indent, #2 entry-spacing
 {\begin{list}{}{%
  \setlength{\itemindent}{0pt}
  \setlength{\leftmargin}{0pt}
  \setlength{\parsep}{0pt}
  % turn on hanging indent if param 1 is 1
  \ifodd #1
   \setlength{\leftmargin}{\cslhangindent}
   \setlength{\itemindent}{-1\cslhangindent}
  \fi
  % set entry spacing
  \setlength{\itemsep}{#2\baselineskip}}}
 {\end{list}}
\usepackage{calc}
\newcommand{\CSLBlock}[1]{\hfill\break\parbox[t]{\linewidth}{\strut\ignorespaces#1\strut}}
\newcommand{\CSLLeftMargin}[1]{\parbox[t]{\csllabelwidth}{\strut#1\strut}}
\newcommand{\CSLRightInline}[1]{\parbox[t]{\linewidth - \csllabelwidth}{\strut#1\strut}}
\newcommand{\CSLIndent}[1]{\hspace{\cslhangindent}#1}
\setlength{\emergencystretch}{3em} % prevent overfull lines
\providecommand{\tightlist}{%
  \setlength{\itemsep}{0pt}\setlength{\parskip}{0pt}}
\usepackage{amsmath}
\usepackage{amssymb}
\usepackage{bookmark}
\IfFileExists{xurl.sty}{\usepackage{xurl}}{} % add URL line breaks if available
\urlstyle{same}
% fallback for those not using the hyperref driver hyperxmp:
\makeatletter
\@ifundefined{xmpquote}{\newcommand{\xmpquote}[1]{#1}}{}
\makeatother
\hypersetup{
  pdftitle={IGNITE: Automatisierte{,} deterministische Thermografie-Pipeline zur Früherkennung pathologischer Entzündungsherde im klinischen Arbeitsumfeld},
  pdfauthor={Jona Noack},
  hidelinks,
  pdfcreator={LaTeX via pandoc}}

\title{IGNITE: Automatisierte, deterministische Thermografie-Pipeline
zur Früherkennung pathologischer Entzündungsherde im klinischen
Arbeitsumfeld}
\author{Jona Noack}
\date{2026-07-23}

\begin{document}
\maketitle

\setstretch{1.5}
\section{Projektüberblick}\label{projektuxfcberblick}

In diesem Projekt für den Wettbewerb \emph{Jugend forscht 2026}
(Fachgebiet Arbeitswelt) wird die Software \textbf{IGNITE} vorgestellt.
Die Arbeit untersucht Möglichkeiten zur Automatisierung der
thermografischen Entzündungserkennung im medizinischen
Behandlungsalltag. Bisher müssen Ärztinnen, Ärzte und Podologen
Wärmebildaufnahmen von Risikopatienten -- beispielsweise beim
diabetischen Fußsyndrom -- manuell am Bildschirm auswerten. Diese
visuelle Sichtprüfung erfordert Zeit (nach eigener Schätzung ca. 3 bis 5
Minuten pro Aufnahme), ist subjektiv und unterliegt tageszeitlichen
Ermüdungsfaktoren des Personals.

Die entwickelte Software nutzt eine 5-stufige mathematische
Bildverarbeitungs-Pipeline, um physiologische Temperaturverläufe zu
filtern, Störabstrahlungen des Raumes zu erodieren und lokale
Hitzespitzen als potenzielle Entzündungsherde zu isolieren. Der
Rechenkern wurde in der Programmiersprache Rust umgesetzt und mittels
Rayon parallelisiert. Auf der nativen Auflösung des verwendeten
Thermosensors (\(160 \times 120\) Pixel) beträgt die gemessene
Auswertungszeit \textbf{1,6 ms} (Median), auf der hochskalierten
JPEG-Ausgabe der Kamera (\(1440 \times 1080\)) dagegen 104 ms. Eine
Kernerkenntnis der Arbeit ist dabei, dass der Rust-Kern \textbf{nicht}
schneller rechnet als die etablierte C++-Bibliothek OpenCV (Faktor 2,4
langsamer); sein Nutzen liegt in der Unabhängigkeit von nativen
Fremdbibliotheken, im deterministischen Verhalten und im geringen
Speicherbedarf (\textless{} 25 MB).

Zur Bewertung wurden \textbf{9 selbst aufgenommene Thermogramme manuell
annotiert} und gegen eine Otsu-Baseline verglichen. IGNITE erreicht
dabei einen Dice-Koeffizienten von \textbf{0,325 ± 0,159} gegenüber
0,004 der Baseline (gepaarter Wilcoxon-Test: \(p = 0{,}0046\)). Bei
getrennter Parameterbestimmung auf einem Tuning-Satz und anschließender
einmaliger Auswertung auf 5 ungesehenen Testbildern steigt der Dice-Wert
auf \textbf{0,513 ± 0,112}. Aufschlussreich ist die Struktur des
Fehlers: Bei einer Precision von 0,99 und einer Sensitivität von 0,21
\textbf{lokalisiert das Verfahren Herde treffsicher, unterschätzt aber
deren Ausdehnung deutlich}. Die zuvor berichteten Werte auf rein
synthetischen Daten (Dice 0,88--0,91) werden in dieser Fassung bewusst
nicht mehr als Wirksamkeitsnachweis geführt, da die simulierten
gaußförmigen Herde exakt der Signalform entsprechen, auf die der Filter
ausgelegt ist. Die Arbeit analysiert die deutlichen Grenzen des
Verfahrens: Da ein deterministischer Filter nicht zwischen biologischen
Infektionen und harmlosen mechanischen Druckstellen (z. B. durch Socken
oder Gehlenken) unterscheiden kann, stellt die Software kein
Medizinprodukt dar, sondern dient als Orientierungshilfe unter
ärztlicher Aufsicht.

\begin{center}\rule{0.5\linewidth}{0.5pt}\end{center}

\section{Inhaltsverzeichnis}\label{inhaltsverzeichnis}

\begin{enumerate}
\def\labelenumi{\arabic{enumi}.}
\tightlist
\item
  \hyperref[1-einleitung-und-problemstellung]{Einleitung und
  Problemstellung}

  \begin{itemize}
  \tightlist
  \item
    1.1 Belastungssituation im medizinischen Behandlungsalltag
  \item
    1.2 Relevanz der Früherkennung beim Diabetischen Fußsyndrom
  \item
    1.3 Zielsetzung und Forschungsfragen
  \end{itemize}
\item
  \hyperref[2-hintergrund-und-theoretische-grundlagen]{Hintergrund und
  theoretische Grundlagen}

  \begin{itemize}
  \tightlist
  \item
    2.1 Medizintechnischer Kontext und podiatrischer Goldstandard
  \item
    2.2 Physikalische Radiometrie und Strahlungsmodell
  \item
    2.3 Kritische Vergleichsmatrix bestehender Auswerteverfahren
  \item
    2.4 Mathematische Ausarbeitung der 5-Stufen-Pipeline
  \end{itemize}
\item
  \hyperref[3-vorgehensweise-materialien-und-methoden]{Vorgehensweise,
  Materialien und Methoden}

  \begin{itemize}
  \tightlist
  \item
    3.1 Analyse und Ergonomie des klinischen Behandlungsablaufs
  \item
    3.2 Software-Architektur und Multi-Backend-Konzept
  \item
    3.3 Schritt-für-Schritt Implementierung in Rust
  \item
    3.4 Ergonomische Benutzeroberfläche und Instant-Splash-UX
  \item
    3.5 Datenschutzkonzept im Arbeitsumfeld
  \item
    3.6 Datenherkunft, Aufnahmegerät und Annotationsverfahren
  \item
    3.7 Selbstständig erbrachter Projektanteil
  \end{itemize}
\item
  \hyperref[4-bildgestuxfctzte-visualisierung-der-pipeline-stufen]{Bildgestützte
  Visualisierung der Pipeline-Stufen}

  \begin{itemize}
  \tightlist
  \item
    4.1 Ausgangsmaterial (Original-Thermogramm)
  \item
    4.2 Körpermaskierung und Distanzkarte (Stufe 2)
  \item
    4.3 Hintergrundkorrektur via Top-Hat-Transformation (Stufe 3)
  \item
    4.4 Statistisches Thresholding und Hotspot-Maske (Stufe 4 \& 5)
  \item
    4.5 Finales diagnostisches Overlay für das Behandlungszimmer
  \item
    4.6 Auswertung synthetischer Krankheits-Szenarien (Regressionstest)
  \end{itemize}
\item
  \hyperref[5-ergebnisse]{Ergebnisse}

  \begin{itemize}
  \tightlist
  \item
    5.1 Laufzeitmessungen und Rechenzeiten
  \item
    5.2 Parameter-Sensitivitätsanalyse des Schwellenwert-Faktors k
  \item
    5.3 Hauptergebnis: Validierung gegen manuell annotierte
    Realaufnahmen
  \item
    5.4 Getrennte Parameterbestimmung: Tuning- und Testsatz
  \item
    5.5 Ablation der geometrischen Filterregeln
  \item
    5.6 Backend-Paritätstests
  \item
    5.7 Ergänzend: synthetische Entzündungsszenarien
  \end{itemize}
\item
  \hyperref[6-ergebnisdiskussion-und-kritische-wuxfcrdigung]{Ergebnisdiskussion
  und Kritische Würdigung}

  \begin{itemize}
  \tightlist
  \item
    6.1 Einordnung der Ergebnisse bezüglich der Arbeitserleichterung
  \item
    6.2 Überprüfung der Hypothesen und Bedeutung von Robust-MAD
  \item
    6.3 Ausführliche Analyse der Nachteile, Grenzen und Störfaktoren
  \end{itemize}
\item
  \hyperref[7-fazit-und-ausblick]{Fazit und Ausblick}

  \begin{itemize}
  \tightlist
  \item
    7.1 Gesamtfazit zur Arbeitswelt-Fragestellung
  \item
    7.2 Zukünftige Erweiterungen für den Praxiseinsatz
  \end{itemize}
\item
  \hyperref[8-literaturverzeichnis]{Literaturverzeichnis}
\item
  \hyperref[9-unterstuxfctzungsleistungen]{Unterstützungsleistungen}
\end{enumerate}

\begin{center}\rule{0.5\linewidth}{0.5pt}\end{center}

\section{1. Einleitung und
Problemstellung}\label{einleitung-und-problemstellung}

\subsection{1.1 Belastungssituation im medizinischen
Behandlungsalltag}\label{belastungssituation-im-medizinischen-behandlungsalltag}

Die demografische Entwicklung und die Zunahme chronischer
Stoffwechselerkrankungen stellen das medizinische Fachpersonal in Praxen
und Kliniken vor erhebliche kapazitäre Herausforderungen (Ring and Ammer
2012). Insbesondere in der Podiatrie, Diabetologie und Dermatologie
müssen täglich zahlreiche Risikopatienten untersucht werden.

Das manuelle Durchmustern von Wärmebildaufnahmen zur Identifikation
pathologischer Hitzespitzen erfordert konzentriertes Absuchen am
Bildschirm, das manuelle Einstellen dynamischer Farbskalen sowie das
Ausmessen kontralateraler Temperaturdifferenzen. Für eine qualifizierte
Erstbeurteilung einer thermografischen Aufnahme ist nach eigener
Zeitmessung an den in dieser Arbeit verwendeten Aufnahmen mit etwa 3 bis
5 Minuten zu rechnen. Es handelt sich ausdrücklich um eine
\textbf{eigene Schätzung} und nicht um einen publizierten Kennwert; eine
belastbare Erhebung des Zeitbedarfs in der Praxis steht noch aus (siehe
Kapitel 6.3). Bei einer Tagesfrequenz von 20 bis 30 Patienten summiert
sich dieser Zusatzaufwand zu einer Arbeitszeit von über 1,5 Stunden, die
für das persönliche Arzt-Patienten-Gespräch fehlt.

Darüber hinaus birgt die rein visuelle Sichtprüfung ein relevantes
Fehlerrisiko: Die wahrgenommene Intensität einer Entzündungszone auf
Farbskalen (z. B. \emph{Rainbow}- oder \emph{Jet}-Colormaps) hängt stark
von der individuellen Farbkontrasteinstellung des Monitors sowie von
Ermüdungsfaktoren des Fachpersonals am Ende einer Schicht ab.

\subsection{1.2 Relevanz der Früherkennung beim Diabetischen
Fußsyndrom}\label{relevanz-der-fruxfcherkennung-beim-diabetischen-fuuxdfsyndrom}

Das diabetische Fußsyndrom (DFS) ist eine Folgeerscheinung der distalen
sensomotorischen Polyneuropathie und der peripheren arteriellen
Verschlusskrankheit (pAVK). Wegen des Verlusts des Protektivempfindens
bleiben biomechanische Überlastungen, Mikrotraumen oder Druckstellen vom
Patienten unbemerkt (Armstrong et al. 2007).

Entzündliche Prozesse im tiefen Gewebe führen durch Hyperämie und
gesteigerte Stoffwechselaktivität zu lokal abgegrenzten
Temperaturerhöhungen der Hautoberfläche. Diese thermischen Anomalien
treten häufig Tage bis Wochen auf, bevor histologische Gewebedefekte
oder ulzeröse Hautdurchbrüche sichtbar werden. Eine verlässliche
Früherkennung ermöglicht frühzeitige Entlastungsmaßnahmen (z. B.
orthopädische Schuhanpassungen) und kann das Risiko von
Unterschenkelamputationen maßgeblich senken (Armstrong et al. 2007).

\subsection{1.3 Zielsetzung und
Forschungsfragen}\label{zielsetzung-und-forschungsfragen}

Ziel dieser Arbeit ist die Konzeption, Implementierung und mathematische
Validierung von \textbf{IGNITE}, einer hochperformanten, lokal
auszuführenden Software zur automatisierten Hotspot-Isolierung. Das
System soll das Fachpersonal durch eine objektive visuelle
Orientierungshilfe entlasten, ohne den klinischen Behandlungsfluss durch
Ladezeiten zu hemmen.

Folgende Forschungsfragen stehen im Zentrum der Untersuchung: *
\textbf{Forschungsfrage 1 (Trefferqualität \& Erklärbarkeit):} Lässt
sich ein deterministischer, mathematisch vollständig nachvollziehbarer
Algorithmus entwickeln, der thermisch auffällige Areale auf \textbf{real
aufgenommenen} Thermogrammen signifikant besser markiert als eine
einfache globale Schwellenwert-Baseline -- ohne auf undurchsichtige
KI-Blackbox-Modelle zurückzugreifen? Als Prüfgröße dient der
Dice-Koeffizient gegenüber einer Handannotation, abgesichert durch einen
gepaarten Signifikanztest. * \textbf{Forschungsfrage 2 (Geschwindigkeit
\& Ergonomie):} Kann eine Rechenzeit von \(< 50\text{ ms}\) erreicht
werden, sodass die Auswertung im Behandlungszimmer verzögerungsfrei
erfolgt -- und welchen Beitrag leistet dabei tatsächlich die
Implementierung in nativem Rust gegenüber der etablierten C++-Bibliothek
OpenCV? Diese Frage wird bewusst als \textbf{Vergleich mit einer
Referenzimplementierung} gestellt und nicht als bloße Bestätigung der
eigenen Umsetzung. * \textbf{Forschungsfrage 3 (Kritische Grenzen):} Wo
liegen die physikalischen und algorithmischen Schwachstellen eines rein
Schwellenwert-basierten Verfahrens im realen Praxisbetrieb?

Vorweggenommen sei, dass die Untersuchung \textbf{zwei der drei
Ausgangserwartungen widerlegt hat}: Der Rust-Kern ist langsamer als
OpenCV (Kapitel 5.1), und die auf synthetischen Daten gemessene
Sensitivität von nahezu 1,00 lässt sich auf Realaufnahmen nicht
reproduzieren (Kapitel 5.3). Beide Befunde werden hier bewusst
berichtet, da sie methodisch aufschlussreicher sind als die ursprünglich
erhofften Bestätigungen.

\begin{center}\rule{0.5\linewidth}{0.5pt}\end{center}

\section{2. Hintergrund und theoretische
Grundlagen}\label{hintergrund-und-theoretische-grundlagen}

\subsection{2.1 Medizintechnischer Kontext und podiatrischer
Goldstandard}\label{medizintechnischer-kontext-und-podiatrischer-goldstandard}

In der medizinischen Thermografie gilt die vergleichende Analyse
symmetrischer Körperareale (kontralaterale Asymmetrie) als
diagnostischer Goldstandard (Armstrong et al. 2007; Ring and Ammer
2012). Da die physiologische Hauttemperatur systemischen Schwankungen
(z. B. Zirkadianer Rhythmus, Raumtemperatur) unterliegt, ist der
absolute Temperaturwert einzelner Pixel nur bedingt aussagekräftig.

Eine kontralaterale Temperaturdifferenz von

\[ \Delta T = |T_{\text{links}} - T_{\text{rechts}}| > 2{,}2\,^\circ\text{C} \]

an anatomisch identischen Messpunkten gilt in der Podiatrie als klinisch
signifikanter Indikator für entzündliche Gewebeprozesse (Armstrong et
al. 1997, 2007).

\pandocbounded{\includegraphics[keepaspectratio,alt={Skizze 4: Kontralaterale Asymmetrie-Analyse}]{images/skizze_asymmetrie_analyse.png}}\\
\emph{Abbildung 1 (Skizze 4): Prinzip der kontralateralen
Asymmetrie-Analyse in der Podiatrie. Das Bild wird an der Bildmitte
getrennt, um die mittleren Oberflächentemperaturen beider Fußsohlen zu
vergleichen. Bei einer Abweichung von
\(\Delta T > 2{,}2\,^\circ\text{C}\) erscheint ein Warnhinweis.}

\begin{center}\rule{0.5\linewidth}{0.5pt}\end{center}

\subsection{2.2 Physikalische Radiometrie und
Strahlungsmodell}\label{physikalische-radiometrie-und-strahlungsmodell}

Die Infrarot-Thermografie basiert auf dem Stefan-Boltzmann-Gesetz
(Stefan 1879; Boltzmann 1884), das die spezifische Ausstrahlung \(M\)
eines schwarzen Körpers in Abhängigkeit von der thermodynamischen
Temperatur \(T\) beschreibt:

\[ M = \sigma \cdot T^4 \]

wobei
\(\sigma \approx 5{,}670374 \times 10^{-8}\,\text{W}\,\text{m}^{-2}\,\text{K}^{-4}\)
die Stefan-Boltzmann-Konstante bezeichnet. Für reale Körper mit dem
spezifischen Emissivitätsgrad \(\epsilon \in (0, 1)\) und unter
Berücksichtigung der von der Umgebung reflektierten Infrarotstrahlung
\(T_{\text{refl}}\) gilt für die vom Sensor erfasste Gewebetemperatur
\(T_{\text{obj}}\) (Jones 1998):

\[ T_{\text{obj}} = \left( \frac{T_{\text{meas}}^4 - (1 - \epsilon) \cdot T_{\text{refl}}^4}{\epsilon} \right)^{1/4} \]

Für menschliches Hautgewebe gilt der messtechnisch bestimmte
Literaturwert \(\epsilon \approx 0{,}98\) (Steketee 1973; Ring and Ammer
2012). Infrarotkameras transformieren das thermische Strahlungsfeld in
eine 8-Bit-Grauwertmatrix \(I(x,y) \in [0, 255]\), deren Intensität
linear mit dem eingestellten Temperaturbereich \([T_{\min}, T_{\max}]\)
korreliert.

\begin{center}\rule{0.5\linewidth}{0.5pt}\end{center}

\subsection{2.3 Kritische Vergleichsmatrix bestehender
Auswerteverfahren}\label{kritische-vergleichsmatrix-bestehender-auswerteverfahren}

{\def\LTcaptype{none} % do not increment counter
\begin{longtable}[]{@{}
  >{\raggedright\arraybackslash}p{(\linewidth - 4\tabcolsep) * \real{0.3333}}
  >{\raggedright\arraybackslash}p{(\linewidth - 4\tabcolsep) * \real{0.3333}}
  >{\raggedright\arraybackslash}p{(\linewidth - 4\tabcolsep) * \real{0.3333}}@{}}
\toprule\noalign{}
\begin{minipage}[b]{\linewidth}\raggedright
Auswerteverfahren
\end{minipage} & \begin{minipage}[b]{\linewidth}\raggedright
Vorteile im Arbeitsalltag
\end{minipage} & \begin{minipage}[b]{\linewidth}\raggedright
Nachteile und Schwachstellen
\end{minipage} \\
\midrule\noalign{}
\endhead
\bottomrule\noalign{}
\endlastfoot
\textbf{Manuelle Sichtprüfung (Arzt/Podologe)} & • Erkennt den
klinischen Gesamtzusammenhang• Kann Narben, Druckstellen \& Wunden
unterscheiden• Benötigt keine Zusatzsoftware & • Zeitaufwendig (3--5
Minuten pro Bild)• Subjektiv und abhängig von Erfahrung/Tagesform• Keine
automatische Dokumentation \\
\textbf{Einfache Schwellenwert-Filter (Otsu)} & • Sehr schnell
(\textless{} 10 ms)• Einfach zu programmieren & • Versagt bei normalen
Körper-Temperaturverläufen• Sehr viele Falsch-Positive an kalten
Rändern \\
\textbf{Deep-Learning KI (z. B. U-Net)} & • Kann komplexe Formen und
Bildmuster erkennen• Hohe Genauigkeit bei gutem Training & •
``Black-Box'': Entscheidungen nicht mathematisch erklärbar• Oft
Cloud-Zwang (DSGVO-Problem im Krankenhaus)• Benötigt leistungsfähige
GPUs \\
\textbf{Mein Ansatz (IGNITE)} & • Schnelle Berechnung (1,6 ms bei
nativer Sensorauflösung, 104 ms bei \(1440\times1080\))• 100 \% lokal \&
DSGVO-konform (kein Cloud-Upload)• Mathematisch nachvollziehbar• Keine
nativen Fremdbibliotheken nötig (\textless{} 25 MB) & • Kann
\textbf{nicht} zwischen Entzündung \& Druckstelle unterscheiden• Feste
Schwellenwerte passen nicht auf jeden Hauttyp• Unterschätzt die
Ausdehnung eines Herdes systematisch (Sensitivität 0,21)• Kein
zertifiziertes Medizinprodukt \\
\end{longtable}
}

\begin{center}\rule{0.5\linewidth}{0.5pt}\end{center}

\subsection{2.4 Mathematische Funktionsweise der
5-Stufen-Pipeline}\label{mathematische-funktionsweise-der-5-stufen-pipeline}

Die Hotspot-Isolierung in IGNITE basiert auf einer 5-stufigen
Bildverarbeitungs-Pipeline (Stiftung Jugend forscht e. V. 2025):

\subsubsection{Stufe 1: Dynamische
Kernel-Skalierung}\label{stufe-1-dynamische-kernel-skalierung}

Um unabhängig von der Sensorauflösung des verwendeten Kamerasystems
(\(160 \times 120\) bis \(1440 \times 1080\) Pixel) identische
geometrische Filtereffekte zu erzielen, skaliert der Radius des
morphologischen Strukturierungselements \(K\) proportional zur minimalen
Bilddimension:

\[ K_{\text{raw}} = \lfloor \min(W, H) \cdot 0{,}05 \rfloor, \quad K_{\text{odd}} = \max(3, K_{\text{raw}} \mid 1) \]

Die bitweise OR-Verknüpfung (\texttt{raw\ \textbar{}\ 1}) erzwingt eine
ungerade Pixelanzahl und garantiert ein eindeutiges mathematisches
Symmetriezentrum.

\subsubsection{Stufe 2: Adaptive Körper-Segmentierung \&
Akren-Erhaltung}\label{stufe-2-adaptive-kuxf6rper-segmentierung-akren-erhaltung}

Zur Abtrennung des kühlen Raumes dient die globale Binarisierung nach
Otsu (Otsu 1979). Bei kontrastarmen Aufnahmen greift ein
Dynamik-Fallback (\(I_{\min} + 0{,}3 \cdot \Delta I\)).

Um isolierte Hintergrundreflexionen (z. B. Falten auf der
Untersuchungsliege) auszuschließen, filtert eine
4-Wege-Zusammenhangskomponentenanalyse alle vom Hauptkörper getrennten
Kleininseln (\(< 2\,\%\) der Gewebefläche) automatisch aus. Da das in
Version 5.0 integrierte 3D-Oberflächenrekonstruktionsmodell den
physikalischen Lambert- und Fresnel-Emissivitätsabfall an Kanten
explizit kompensiert, erfordert die Maskenerosion keinen aggressiven
Randverschnitt mehr: Der relative Distanzfaktor wird auf
\(f_{\text{dist}} \le 0{,}005\) (\(1\dots 2\,\text{px}\)) minimiert.
Schmale distale Extremitäten (wie Zehen und Finger) bleiben so
vollständig in der diagnostischen Maske erhalten:

\[ \text{Mask}_{\text{eroded}}(x,y) = \begin{cases} 255, & \text{falls } D(x,y) \ge f_{\text{dist}} \cdot \max_{x',y'} D(x',y') \\ 0, & \text{sonst} \end{cases} \]

\subsubsection{Stufe 3: Morphologische
Top-Hat-Transformation}\label{stufe-3-morphologische-top-hat-transformation}

Zur Elimination physiologischer Helligkeitsverläufe (z. B. der
natürlichen Gewewärme im Fußgewölbe) wird die morphologische
Top-Hat-Transformation verwendet:

\[ \text{TopHat}(I) = I - \gamma_K(I) = I - ((I \ominus K) \oplus K) \]

wobei \(\gamma_K(I)\) das morphologische Opening von \(I\) mit dem
Kernel \(K\) bezeichnet. Im Rust-Core ist die 2D-Operation in zwei
sequentielle 1D-Durchläufe (horizontal und vertikal) nach Lemire (Lemire
2006) zerlegt. Die Komplexität pro Pixel sinkt dadurch von \(O(K^2)\)
auf \(O(K)\).

\pandocbounded{\includegraphics[keepaspectratio,alt={Skizze 1: Prinzip der morphologischen Top-Hat-Transformation}]{images/skizze_tophat_prinzip.png}}\\
\emph{Abbildung 2 (Skizze 1): Das Prinzip der morphologischen
Top-Hat-Transformation im 1D-Temperaturprofil. Das morphologische
Opening glättet großflächige Verläufe. Die Differenz isoliert scharfe
lokale Hitzespitzen oberhalb des Schwellenwerts.}

\begin{center}\rule{0.5\linewidth}{0.5pt}\end{center}

\subsubsection{Stufe 4: Statistisches Outlier-Thresholding (Gauß
vs.~Robust-MAD)}\label{stufe-4-statistisches-outlier-thresholding-gauuxdf-vs.-robust-mad}

Zur Segmentierung signifikanter Hitzespitzen werden zwei statistische
Verfahren unterstützt: 1. \textbf{Gaussian Thresholding:}

\[ T_{\text{Gauß}} = \mu_{\text{diff}} + k \cdot \sigma_{\text{diff}} \]

\begin{enumerate}
\def\labelenumi{\arabic{enumi}.}
\setcounter{enumi}{1}
\tightlist
\item
  \textbf{Robustes MAD-Thresholding:} Bei bimodalen
  Temperaturverteilungen (z. B. stark unterkühlten Zehen) verzerrt der
  kalte Pol den Mittelwert \(\mu\). IGNITE berechnet in diesem Fall die
  robusten Kennzahlen Median (\(\tilde{\mu}\)) und Median Absolute
  Deviation (\(\text{MAD}\)):
\end{enumerate}

\[ \text{MAD} = \text{median}(|X - \tilde{\mu}|), \quad \hat{\sigma}_{\text{MAD}} = 1{,}4826 \cdot \text{MAD} \]

\[ T_{\text{MAD}} = \tilde{\mu} + k \cdot \hat{\sigma}_{\text{MAD}} \]

Um zu verhindern, dass bei bimodalen Verteilungen mit stark unterkühlten
Extremitäten (z. B. Zehen oder Fingerkuppen bei 24--30 °C) diese kühlen
Regionen durch einen starren globalen Medianschnitt
(\(I \ge \tilde{\mu}\)) vollständig verworfen werden, implementiert das
System eine adaptive physiologische Gewebeschwelle:

\[ T_{\text{floor}} = \min(\max(\tilde{\mu} \cdot 0{,}55, 45), 80) \]

Dadurch werden echte Entzündungsherde auf kühlen Akren verlässlich
detektiert, während Raumluftartefakte (\textless{} 20 °C) sicher
eliminiert bleiben.

\pandocbounded{\includegraphics[keepaspectratio,alt={Skizze 2: Gauß vs.~Robust-MAD bei bimodaler Verteilung}]{images/skizze_gauss_vs_mad.png}}\\
\emph{Abbildung 3 (Skizze 2): Vergleichende Skizze der statistischen
Schwellenwerte bei einer bimodalen Gewebeverteilung (kalte Zehen). Der
Gauß-Mittelwert verschiebt sich nach links und verzerrt die Schwelle,
während das Median/MAD-Verfahren stabil bleibt.}

\begin{center}\rule{0.5\linewidth}{0.5pt}\end{center}

\subsubsection{Stufe 5: Geometrische Rauschfilterung \&
Asymmetrie}\label{stufe-5-geometrische-rauschfilterung-asymmetrie}

Zusammenhängende Pixelgruppen werden über eine Zusammenhangsanalyse in
4er-Nachbarschaft bestimmt. Ein Cluster wird \textbf{beibehalten}, wenn
seine Fläche mindestens \(0{,}05\,\%\) der Körperoberfläche beträgt und
seine Form-Circularity \(C\) den Schwellenwert erreicht:

\[ C = \frac{4\pi \cdot A}{P^2} \ge 0{,}08 \]

Der Wert \(0{,}08\) entspricht der Implementierung
(\texttt{DEFAULT\_MIN\_CIRCULARITY} in \texttt{config.py}). Er ist
bewusst niedrig gewählt: Ein exakter Kreis erreicht \(C = 1\), reale
Entzündungsherde sind jedoch unregelmäßig berandet, und durch die
pixelweise Umfangsberechnung wird \(P\) bei kleinen Objekten
systematisch überschätzt, was \(C\) zusätzlich verringert. Der Filter
entfernt daher gezielt langgestreckte, fadenförmige Strukturen (etwa
Kanten- und Kompressionsartefakte), ohne kompakte Herde zu verwerfen.

In Version 5.0 wurde die vormalige starre 65-\%-Grenze
(\texttt{AnatomicalCutoffY}) vollständig überwunden: Das System arbeitet
\textbf{strikt anatomie-agnostisch} und eignet sich gleichermaßen für
Hände/Finger, Knie, Rücken/Wirbelsäule und Füße. Durch die
biophysikalische Randprüfung (Erhalt von Herden mit thermischem
Diffusionsgradienten \(G_{\text{edge}} \ge 3{,}0\) oder Halo
\(\Delta T_{\text{halo}} \ge 5{,}0\)) werden echte Entzündungsherde an
Zehen- und Fingerkuppen zuverlässig erfasst, während flache optische
Randleckartefakte zuverlässig verworfen werden.

\begin{center}\rule{0.5\linewidth}{0.5pt}\end{center}

\section{3. Vorgehensweise, Materialien und
Methoden}\label{vorgehensweise-materialien-und-methoden}

\subsection{3.1 Analyse des klinischen
Behandlungsablaufs}\label{analyse-des-klinischen-behandlungsablaufs}

Der gedachte Ablauf im Behandlungszimmer gliedert sich wie folgt:

\pandocbounded{\includegraphics[keepaspectratio,alt={Skizze 3: Integration in den Praxis-Workflow}]{images/skizze_praxis_workflow.png}}\\
\emph{Abbildung 4 (Skizze 3): Schema der Integration von IGNITE in den
täglichen Praxisablauf im Behandlungszimmer.}

Die Software dient dabei als Werkzeug für Schritt 3. Die eigentliche
Diagnose in Schritt 4 bleibt immer beim Fachpersonal.

\subsection{3.2 Software-Architektur und
Speicherverwaltung}\label{software-architektur-und-speicherverwaltung}

Die Software ist modular aufgebaut. Um maximale Geschwindigkeit zu
erreichen und Speicherzugriffe zu minimieren, werden die Bilddaten über
die C-ABI ohne Kopiervorgänge (Zero-Copy) direkt aus dem Python-Speicher
in den Rust-Core übergeben:

\pandocbounded{\includegraphics[keepaspectratio,alt={Skizze 5: Rust FFI \& Speicherarchitektur}]{images/skizze_rust_ffi_architektur.png}}\\
\emph{Abbildung 5 (Skizze 5): Speicherarchitektur und FFI-Anbindung
zwischen Python (NumPy) und dem Rust-Core (\texttt{ignite\_core}). Die
Bilddaten werden ohne Kopieren via Zeiger übergeben und in Rust mit
Rayon auf CPU-Kerne verteilt.}

\begin{center}\rule{0.5\linewidth}{0.5pt}\end{center}

\subsection{3.3 Schritt-für-Schritt Implementierung in
Rust}\label{schritt-fuxfcr-schritt-implementierung-in-rust}

\begin{enumerate}
\def\labelenumi{\arabic{enumi}.}
\tightlist
\item
  \textbf{Python-Prototyp:} Erste Tests zeigten, dass OpenCV in Python
  bei großen Bildern 80 bis 210 ms benötigte.
\item
  \textbf{Rust-Umsetzung:} Durch die Umschreibung in Rust mit den Crates
  \texttt{ndarray} und \texttt{imageproc} konnte die Zeit gesenkt
  werden.
\item
  \textbf{Parallelisierung:} Mit Rayon (\texttt{par\_iter()}) werden die
  Schleifen auf alle verfügbaren CPU-Kerne verteilt.
\end{enumerate}

\subsection{3.4 Benutzeroberfläche und
Ladeoptimierung}\label{benutzeroberfluxe4che-und-ladeoptimierung}

Um lange Startzeiten durch schwere Bibliotheken zu vermeiden, öffnet
\texttt{main.py} beim Aufruf in unter 50 ms einen leichten
Tkinter-Splash-Screen. Während der Anwender die Rückmeldung sieht, lädt
ein Hintergrund-Thread die benötigten Module.

\subsection{3.5 Datenschutzkonzept}\label{datenschutzkonzept}

\begin{enumerate}
\def\labelenumi{\arabic{enumi}.}
\tightlist
\item
  \textbf{In-Memory:} Es werden keine Bilddaten auf externe Server
  übertragen.
\item
  \textbf{SHA-256 Pseudonymisierung:} Patientennamen werden mit Salt
  gehasht (\texttt{ANON-\textless{}hash\textgreater{}}).
\end{enumerate}

\subsection{3.6 Datenherkunft, Aufnahmegerät und
Annotationsverfahren}\label{datenherkunft-aufnahmegeruxe4t-und-annotationsverfahren}

\subsubsection{Herkunft der Aufnahmen}\label{herkunft-der-aufnahmen}

Alle 21 in dieser Arbeit ausgewerteten Thermogramme wurden \textbf{von
mir selbst aufgenommen}. Als Aufnahmepersonen dienten
\textbf{Familienmitglieder und Bekannte}, die vor der Aufnahme über
Zweck und Verwendung der Bilder informiert wurden und ihr
\textbf{Einverständnis} erteilt haben. Bei minderjährigen Personen lag
zusätzlich das Einverständnis der Erziehungsberechtigten vor.

Um Fehlinterpretationen vorzubeugen, ist ausdrücklich festzuhalten:

\begin{itemize}
\tightlist
\item
  Es handelt sich \textbf{nicht um klinische Patientendaten}. Ein
  früherer Wortlaut dieser Arbeit sprach von „21 realen klinischen
  Testaufnahmen''; diese Formulierung war irreführend und wurde
  korrigiert.
\item
  Die aufgenommenen Personen sind \textbf{nicht diagnostiziert
  erkrankt}. Es liegen daher keine ärztlich bestätigten Entzündungsherde
  vor; die Annotationen markieren \emph{thermisch auffällige Regionen},
  nicht \emph{gesicherte Pathologien}.
\item
  Es wurde \textbf{kein Ethikvotum} eingeholt, da keine Studie an
  Patientinnen und Patienten durchgeführt wurde und keine
  Gesundheitsdaten Dritter aus einer Behandlung stammen.
\item
  Die Bilddateien enthalten keine Namen; eine Zuordnung zu Personen ist
  über die Dateinamen (\texttt{bild\ (1)} \ldots{} \texttt{bild\ (21)})
  nicht möglich.
\end{itemize}

\subsubsection{Aufnahmegerät und dessen
Konsequenzen}\label{aufnahmegeruxe4t-und-dessen-konsequenzen}

Verwendet wurde eine \textbf{FLIR ONE bzw. FLIR ONE Pro} -- eine an ein
Smartphone ansteckbare Thermografiekamera. Für die Bewertung der
Ergebnisse ist eine Eigenschaft dieses Geräts von zentraler Bedeutung:

\begin{quote}
Der \textbf{thermische Sensor} der FLIR ONE Pro besitzt eine Auflösung
von lediglich \textbf{\(160 \times 120\) Pixeln}. Die ausgegebene
JPEG-Datei mit \(1440 \times 1080\) Pixeln entsteht durch Hochskalierung
und Überlagerung mit den Kanten des sichtbaren Lichtbildes (\emph{MSX}).
Sie enthält \textbf{keine zusätzliche thermische Information}.
\end{quote}

Daraus folgen drei wesentliche Einschränkungen:

\begin{enumerate}
\def\labelenumi{\arabic{enumi}.}
\tightlist
\item
  \textbf{Scheinauflösung.} Ein im \(1440 \times 1080\)-Bild fein
  umrandeter Hotspot beruht real auf wenigen Sensorpixeln. Die
  Genauigkeit der Konturbestimmung ist entsprechend begrenzt -- dies
  erklärt einen Teil der in Kapitel 5.3 gemessenen Abweichung zwischen
  Algorithmus- und Handannotation.
\item
  \textbf{Verlust der Radiometrie.} Die exportierten JPEG-Dateien
  enthalten \textbf{keine Absoluttemperaturen}, sondern nur
  8-Bit-Grauwerte einer geräteseitig und je Bild dynamisch skalierten
  Farbpalette. Die in Kapitel 2.2 hergeleitete Umrechnung nach
  Stefan-Boltzmann konnte auf diesen Daten daher \textbf{nicht
  angewendet} werden. Insbesondere ist das podiatrische Kriterium einer
  kontralateralen Differenz von \(\Delta T > 2{,}2\,\)K (Armstrong et
  al. 1997) auf diesem Material \textbf{nicht überprüfbar}; die Pipeline
  arbeitet ausschließlich auf relativen Intensitäten.
\item
  \textbf{Kompressionsartefakte.} Die JPEG-Kompression erzeugt
  Blockartefakte, die der Top-Hat-Filter prinzipiell als lokale Maxima
  aufgreifen kann.
\end{enumerate}

\subsubsection{Aufnahmebedingungen}\label{aufnahmebedingungen}

Die Aufnahmen erfolgten in Innenräumen bei üblicher Zimmertemperatur,
aus etwa gleichem Abstand und in gleicher Ausrichtung (Füße mittig im
Bild). Eine standardisierte Akklimatisierungsphase, wie sie das
\emph{Glamorgan-Protokoll} für medizinische Thermografie vorsieht, wurde
\textbf{nicht} eingehalten. Raumtemperatur und Luftbewegung wurden nicht
protokolliert. Die einheitliche Aufnahmegeometrie ist zugleich der Grund
dafür, dass die geometrischen Filterregeln aus Kapitel 5.5 auf diesem
Datensatz unschädlich sind -- auf anders aufgenommenem Material wäre das
nicht zu erwarten.

\subsubsection{Annotationsverfahren}\label{annotationsverfahren}

Die Referenzmasken wurden mit dem eigens entwickelten Werkzeug
\texttt{benchmark\_annotator.py} per Pinselwerkzeug erstellt und als
binäre PNG-Masken abgelegt. Wichtige Einschränkungen:

\begin{itemize}
\tightlist
\item
  Annotiert wurden \textbf{9 der 21 Aufnahmen}; für die übrigen 12 liegt
  keine Referenz vor.
\item
  Die Annotation erfolgte durch \textbf{eine einzige, medizinisch nicht
  ausgebildete Person} (den Autor). Es gibt daher \textbf{weder eine
  Intra- noch eine Inter-Rater-Übereinstimmung} (z. B. Cohens
  \(\kappa\)), und die Referenz ist \textbf{kein diagnostischer
  Goldstandard}.
\item
  Die Annotation erfolgte in Kenntnis des Projektziels, eine Verblindung
  fand nicht statt.
\end{itemize}

Diese Punkte begrenzen die Aussagekraft aller in Kapitel 5.3 und 5.4
berichteten Werte erheblich und werden in Kapitel 6.3 erneut
aufgegriffen.

\subsection{3.7 Selbstständig erbrachter
Projektanteil}\label{selbststuxe4ndig-erbrachter-projektanteil}

Die mathematische Ausarbeitung, die Programmierung in Rust und Python,
das Erstellen der Benutzeroberfläche sowie die Durchführung aller Tests
wurden zu 100 \% eigenständig von mir durchgeführt.

\begin{center}\rule{0.5\linewidth}{0.5pt}\end{center}

\section{4. Bildgestützte Visualisierung der
Pipeline-Stufen}\label{bildgestuxfctzte-visualisierung-der-pipeline-stufen}

Die folgenden Abbildungen zeigen die Zwischenschritte der Pipeline an
einem echten Testbild:

\subsubsection{4.1 Ausgangsmaterial
(Original-Thermogramm)}\label{ausgangsmaterial-original-thermogramm}

\pandocbounded{\includegraphics[keepaspectratio,alt={Originales Wärmebild in Grauwert und Jet-Colormap}]{images/1_original_thermal_jet.png}}\\
\emph{Abbildung 6: Originales thermografisches Wärmebild in
Jet-Colormap.}

\begin{center}\rule{0.5\linewidth}{0.5pt}\end{center}

\subsubsection{4.2 Körpermaskierung und Distanzkarte (Stufe
2)}\label{kuxf6rpermaskierung-und-distanzkarte-stufe-2}

\pandocbounded{\includegraphics[keepaspectratio,alt={Körpermaske und Chamfer-L2 Distanzkarte}]{images/2_distance_transform.png}}\\
\emph{Abbildung 7: Chamfer-L2 Distanzkarte zur Abtrennung unscharfer
Ränder.}

\begin{center}\rule{0.5\linewidth}{0.5pt}\end{center}

\subsubsection{4.3 Hintergrundkorrektur via Top-Hat-Transformation
(Stufe
3)}\label{hintergrundkorrektur-via-top-hat-transformation-stufe-3}

\pandocbounded{\includegraphics[keepaspectratio,alt={Top-Hat Differenzbild}]{images/3_tophat_difference.png}}\\
\emph{Abbildung 8: Ergebnis der Top-Hat-Transformation (isoliert lokale
Temperaturunterschiede).}

\begin{center}\rule{0.5\linewidth}{0.5pt}\end{center}

\subsubsection{4.4 Statistisches Thresholding und Hotspot-Maske (Stufe 4
\& 5)}\label{statistisches-thresholding-und-hotspot-maske-stufe-4-5}

\pandocbounded{\includegraphics[keepaspectratio,alt={Binäre Hotspot-Maske}]{images/4_hotspot_mask.png}}\\
\emph{Abbildung 9: Binäre Hotspot-Maske nach Schwellenwertentscheidung
und Rauschfilterung.}

\begin{center}\rule{0.5\linewidth}{0.5pt}\end{center}

\subsubsection{4.5 Finales diagnostisches
Overlay}\label{finales-diagnostisches-overlay}

\pandocbounded{\includegraphics[keepaspectratio,alt={Finales Hotspot Overlay}]{images/5_final_overlay_jet.png}}\\
\emph{Abbildung 10: Visuelle Orientierungshilfe mit roter
Hotspot-Markierung.}

\begin{center}\rule{0.5\linewidth}{0.5pt}\end{center}

\subsubsection{4.6 Auswertung synthetischer
Szenarien}\label{auswertung-synthetischer-szenarien}

\pandocbounded{\includegraphics[keepaspectratio,alt={Synthetisches Szenario Diabetischer Fuß}]{images/synthetic_diabetic_ulcer.png}}\\
\emph{Abbildung 11: Auswertung eines simulierten diabetischen
Fußgeschwürs.}

\begin{center}\rule{0.5\linewidth}{0.5pt}\end{center}

\section{5. Ergebnisse}\label{ergebnisse}

\subsection{5.1 Laufzeitmessungen und
Rechenzeiten}\label{laufzeitmessungen-und-rechenzeiten}

Alle Zahlen dieses Kapitels werden von
\texttt{scripts/run\_validation.py} reproduzierbar erzeugt (fester
Zufalls-Seed, protokollierte Hardware). Der vollständige Ergebnisbericht
ist als \texttt{docs/validation/validation\_report.json} dem Projekt
beigelegt. Gemessen wurde auf einem x86\_64-System mit \textbf{4
CPU-Kernen}, 60 Wiederholungen nach 10 Aufwärmläufen. Angegeben sind
Minimum und Median, da Einzelmessungen unter einem
Desktop-Betriebssystem stark streuen.

{\def\LTcaptype{none} % do not increment counter
\begin{longtable}[]{@{}
  >{\raggedright\arraybackslash}p{(\linewidth - 4\tabcolsep) * \real{0.2857}}
  >{\centering\arraybackslash}p{(\linewidth - 4\tabcolsep) * \real{0.3571}}
  >{\centering\arraybackslash}p{(\linewidth - 4\tabcolsep) * \real{0.3571}}@{}}
\toprule\noalign{}
\begin{minipage}[b]{\linewidth}\raggedright
Backend
\end{minipage} & \begin{minipage}[b]{\linewidth}\centering
\(400 \times 400\) (min / median)
\end{minipage} & \begin{minipage}[b]{\linewidth}\centering
\(1440 \times 1080\) (min / median)
\end{minipage} \\
\midrule\noalign{}
\endhead
\bottomrule\noalign{}
\endlastfoot
\textbf{Rust-Core (CPU, rayon)} & 7,9 / 9,1 ms & 86,4 / 104,3 ms \\
\textbf{Python + OpenCV (CPU)} & 3,1 / 3,3 ms & 40,6 / 43,9 ms \\
\textbf{PyTorch (CUDA)} & nicht validierbar & nicht validierbar \\
\end{longtable}
}

\textbf{Dieses Ergebnis widerlegt meine ursprüngliche Erwartung und ist
das wichtigste methodische Resultat der Laufzeituntersuchung.} In einer
früheren Messung erschien der Rust-Core rund dreimal schneller als der
Python-Pfad. Diese Messung war jedoch \textbf{nicht fair}: Der
Python-Fallback berechnete damals eine
\emph{Multi-Scale}-Top-Hat-Transformation über drei Skalen zuzüglich
einer Vorglättung, während der Rust-Core nur \textbf{eine} Skala
rechnete. Verglichen wurden also unterschiedliche Algorithmen. Nachdem
beide Backends auf denselben einskaligen Algorithmus angeglichen wurden
(siehe Kapitel 5.5), zeigt sich:

\begin{quote}
Die SIMD- und multithread-optimierte C++-Morphologie von OpenCV ist auf
dieser Hardware etwa \textbf{2,4-mal schneller} als meine
handgeschriebene Rust-Implementierung.
\end{quote}

Das ist kein Fehler der Rust-Umsetzung, sondern erwartbar: OpenCV ist
eine seit über zwei Jahrzehnten auf Vektorbefehlssätze hin optimierte
Bibliothek. Der Nutzen des Rust-Cores liegt daher \textbf{nicht in der
Rohgeschwindigkeit}, sondern in:

\begin{itemize}
\tightlist
\item
  \textbf{Unabhängigkeit von nativen Fremdbibliotheken} -- der Installer
  benötigt keine OpenCV-Laufzeitkomponenten (\textless{} 25 MB statt
  \textgreater{} 200 MB),
\item
  \textbf{deterministischem, plattformgleichem Verhalten} ohne
  Abhängigkeit von der jeweils installierten OpenCV-Version,
\item
  \textbf{Speichersicherheit} und geringem Arbeitsspeicherbedarf
  (\textless{} 25 MB).
\end{itemize}

\subsubsection{Profilierung der
Pipeline-Stufen}\label{profilierung-der-pipeline-stufen}

Mit gesetzter Umgebungsvariable \texttt{IGNITE\_DEBUG} gibt der
Rust-Core Stufenzeiten aus (\(1440 \times 1080\), Minimum aus 10
Läufen):

{\def\LTcaptype{none} % do not increment counter
\begin{longtable}[]{@{}lc@{}}
\toprule\noalign{}
Stufe & Zeit \\
\midrule\noalign{}
\endhead
\bottomrule\noalign{}
\endlastfoot
Körpermaske (Otsu + Closing + Distanztransformation) & 39,0 ms \\
Top-Hat-Transformation & 26,8 ms \\
Statistisches Thresholding & 2,5 ms \\
Geometriefilter & 5,7 ms \\
\end{longtable}
}

Die Profilierung führte zu einer gezielten Optimierung: Das statistische
Thresholding materialisierte zuvor zwei Vektoren mit je rund 500 000
\texttt{f64}-Werten. Durch eine einzige parallele Reduktion über Summen
(ohne Zwischenspeicherung) sank diese Stufe von 17,9 ms auf 2,5 ms --
eine Beschleunigung um Faktor 7.

\subsubsection{Auflösungsabhängigkeit und praktische
Konsequenz}\label{aufluxf6sungsabhuxe4ngigkeit-und-praktische-konsequenz}

Die FLIR ONE Pro besitzt einen \textbf{thermischen Sensor von nur
\(160 \times 120\) Pixeln}; die ausgegebene JPEG-Datei mit
\(1440 \times 1080\) Pixeln ist eine hochskalierte Darstellung (vgl.
Kapitel 3.6). Die Verarbeitung der hochskalierten Datei kostet daher
Rechenzeit, ohne zusätzliche thermische Information zu liefern:

{\def\LTcaptype{none} % do not increment counter
\begin{longtable}[]{@{}lc@{}}
\toprule\noalign{}
Verarbeitungsauflösung & Rust-Core (min / median) \\
\midrule\noalign{}
\endhead
\bottomrule\noalign{}
\endlastfoot
\(160 \times 120\) (native Sensorauflösung) & \textbf{1,2 / 1,6 ms} \\
\(320 \times 240\) & 3,8 / 4,7 ms \\
\(640 \times 480\) & 17,5 / 20,1 ms \\
\(1440 \times 1080\) (JPEG-Ausgabegröße) & 86,4 / 104,3 ms \\
\end{longtable}
}

\textbf{Forschungsfrage 2 ist damit differenziert zu beantworten:} Auf
der hochskalierten JPEG-Auflösung wird die ursprünglich angestrebte
Grenze von 50 ms \textbf{nicht} eingehalten (86--104 ms). Auf der
tatsächlichen Sensorauflösung liegt die Auswertung dagegen mit 1,6 ms um
mehr als das Dreißigfache unter dieser Grenze. Für den Praxiseinsatz
folgt daraus die konkrete Empfehlung, vor der Analyse auf die native
Sensorauflösung zu reduzieren, statt auf Interpolationsartefakten zu
rechnen.

\subsection{5.2 Parameter-Sensitivitätsanalyse des Schwellenwert-Faktors
k}\label{parameter-sensitivituxe4tsanalyse-des-schwellenwert-faktors-k}

Um zu untersuchen, wie stabil der Schwellenwert-Faktor \(k\) auf das
Filterergebnis reagiert, habe ich eine Sensitivitätsanalyse
durchgeführt:

\pandocbounded{\includegraphics[keepaspectratio,alt={Diagramm 2: Parameter-Sensitivitätsanalyse}]{images/diagramm_parameter_sensitivitaet.png}}\\
\emph{Abbildung 13 (Diagramm 2): Sensitivitätsanalyse des Faktors \(k\)
auf \textbf{synthetischen} Daten.}

\textbf{Wichtige Einschränkung:} Diese Analyse beruht auf synthetischen
Szenarien und legte den Bereich \(k \in [2{,}5;\ 3{,}5]\) nahe. Die
Auswertung gegen echte Handannotationen (Kapitel 5.4) widerspricht dem
deutlich: Dort liegt das Optimum bei \(k = 1{,}25\), während
\(k = 3{,}0\) den Dice-Wert um rund ein Drittel verschlechtert. Die
synthetische Sensitivitätsanalyse ist daher \textbf{nicht auf reale
Aufnahmen übertragbar} -- ein weiterer Beleg dafür, dass simulierte
Herde die Eigenschaften echter Thermogramme nur unzureichend abbilden.
Maßgeblich für die Parameterwahl ist ausschließlich Kapitel 5.4.

\begin{center}\rule{0.5\linewidth}{0.5pt}\end{center}

\subsection{5.3 Hauptergebnis: Validierung gegen manuell annotierte
Realaufnahmen}\label{hauptergebnis-validierung-gegen-manuell-annotierte-realaufnahmen}

Dies ist das zentrale Ergebnis der Arbeit. Ausgewertet wurden die
\textbf{9 von 21 Aufnahmen, für die eine manuell erstellte
Ground-Truth-Maske vorliegt} (Annotationswerkzeug:
\texttt{benchmark\_annotator.py}). Verglichen wird gegen eine
\textbf{Otsu-Baseline} -- eine einfache globale Schwellenwertbildung
ohne Top-Hat und ohne Geometriefilter.

Angegeben sind Mittelwert ± Standardabweichung sowie das
95-\%-Bootstrap-Konfidenz\-intervall (10 000 Ziehungen, Seed 20260726)
bei Standardparametern (\(k = 3{,}0\)):

{\def\LTcaptype{none} % do not increment counter
\begin{longtable}[]{@{}lccc@{}}
\toprule\noalign{}
Metrik & IGNITE & 95 \%-KI & Otsu-Baseline \\
\midrule\noalign{}
\endhead
\bottomrule\noalign{}
\endlastfoot
\textbf{Dice} & \textbf{0,325 ± 0,159} & 0,240 -- 0,429 & 0,004 ±
0,011 \\
\textbf{IoU} & 0,205 ± 0,132 & 0,137 -- 0,292 & 0,002 ± 0,006 \\
\textbf{Sensitivität} & 0,206 ± 0,133 & 0,138 -- 0,293 & 0,005 ±
0,011 \\
\textbf{Spezifität} & 1,000 ± 0,000 & 1,000 -- 1,000 & 0,982 ± 0,015 \\
\textbf{Precision} & \textbf{0,991 ± 0,022} & 0,976 -- 1,000 & -- \\
\textbf{MCC} & 0,434 ± 0,126 & 0,367 -- 0,516 & -- \\
\end{longtable}
}

Ein gepaarter \textbf{Wilcoxon-Vorzeichen-Rangtest} über die neun
Bildpaare ergibt einen signifikanten Vorteil von IGNITE gegenüber der
Otsu-Baseline (\(W = 45{,}0\); \(p = 0{,}0046\); mittlere Dice-Differenz
\(+0{,}321\)).

\subsubsection{Interpretation: hohe Präzision bei geringer
Flächenabdeckung}\label{interpretation-hohe-pruxe4zision-bei-geringer-fluxe4chenabdeckung}

Das Ergebnismuster ist eindeutig und diagnostisch aufschlussreich: Bei
einer \textbf{Precision von 0,991} liegt praktisch \textbf{jeder} von
IGNITE markierte Pixel innerhalb der ärztlich annotierten Region --
Fehlalarme treten also kaum auf. Gleichzeitig liegt die
\textbf{Sensitivität bei nur 0,206}, das heißt IGNITE markiert im Mittel
lediglich rund ein Fünftel der annotierten Fläche.

IGNITE \textbf{lokalisiert thermisch auffällige Regionen also
treffsicher, unterschätzt aber systematisch deren Ausdehnung.}
Ursächlich ist der Top-Hat-Filter: Er reagiert auf den steilen
Temperaturgradienten im Zentrum eines Herdes, während die flach
auslaufenden Randzonen unter dem statistischen Schwellenwert bleiben.
Für die angestrebte Rolle als \emph{Orientierungshilfe} -- das Auge des
Fachpersonals auf die richtige Stelle zu lenken -- ist dieses Verhalten
günstiger als der umgekehrte Fall (viele Fehlalarme). Für eine
Flächenvermessung einer Wunde ist das Verfahren dagegen \textbf{nicht}
geeignet.

\begin{center}\rule{0.5\linewidth}{0.5pt}\end{center}

\subsection{5.4 Getrennte Parameterbestimmung: Tuning- und
Testsatz}\label{getrennte-parameterbestimmung-tuning--und-testsatz}

In einer früheren Fassung dieser Arbeit wurde der Schwellenwertfaktor
\(k\) auf denselben Bildern bestimmt, auf denen anschließend die Güte
berichtet wurde. Das ist methodisch unzulässig (\emph{Data Leakage}) und
führt zu systematisch zu optimistischen Werten. Die Auswertung wurde
daher auf eine \textbf{strikte Trennung} umgestellt:

\begin{itemize}
\tightlist
\item
  Die 9 annotierten Bilder werden mit festem Seed (20260726) zufällig
  geteilt in \textbf{4 Tuning-Bilder} und \textbf{5 Testbilder}.
\item
  Auf dem Tuning-Satz wird \(k\) über ein Raster
  \(k \in [0{,}5;\ 4{,}5]\) gewählt.
\item
  Der Testsatz wird \textbf{genau einmal} mit dem gewählten \(k\)
  ausgewertet.
\end{itemize}

\textbf{Dice auf dem Tuning-Satz in Abhängigkeit von \(k\):}

{\def\LTcaptype{none} % do not increment counter
\begin{longtable}[]{@{}
  >{\raggedright\arraybackslash}p{(\linewidth - 22\tabcolsep) * \real{0.1081}}
  >{\centering\arraybackslash}p{(\linewidth - 22\tabcolsep) * \real{0.0811}}
  >{\centering\arraybackslash}p{(\linewidth - 22\tabcolsep) * \real{0.0811}}
  >{\centering\arraybackslash}p{(\linewidth - 22\tabcolsep) * \real{0.0811}}
  >{\centering\arraybackslash}p{(\linewidth - 22\tabcolsep) * \real{0.0811}}
  >{\centering\arraybackslash}p{(\linewidth - 22\tabcolsep) * \real{0.0811}}
  >{\centering\arraybackslash}p{(\linewidth - 22\tabcolsep) * \real{0.0811}}
  >{\centering\arraybackslash}p{(\linewidth - 22\tabcolsep) * \real{0.0811}}
  >{\centering\arraybackslash}p{(\linewidth - 22\tabcolsep) * \real{0.0811}}
  >{\centering\arraybackslash}p{(\linewidth - 22\tabcolsep) * \real{0.0811}}
  >{\centering\arraybackslash}p{(\linewidth - 22\tabcolsep) * \real{0.0811}}
  >{\centering\arraybackslash}p{(\linewidth - 22\tabcolsep) * \real{0.0811}}@{}}
\toprule\noalign{}
\begin{minipage}[b]{\linewidth}\raggedright
\(k\)
\end{minipage} & \begin{minipage}[b]{\linewidth}\centering
0,5
\end{minipage} & \begin{minipage}[b]{\linewidth}\centering
0,75
\end{minipage} & \begin{minipage}[b]{\linewidth}\centering
1,0
\end{minipage} & \begin{minipage}[b]{\linewidth}\centering
\textbf{1,25}
\end{minipage} & \begin{minipage}[b]{\linewidth}\centering
1,5
\end{minipage} & \begin{minipage}[b]{\linewidth}\centering
2,0
\end{minipage} & \begin{minipage}[b]{\linewidth}\centering
2,5
\end{minipage} & \begin{minipage}[b]{\linewidth}\centering
3,0
\end{minipage} & \begin{minipage}[b]{\linewidth}\centering
3,5
\end{minipage} & \begin{minipage}[b]{\linewidth}\centering
4,0
\end{minipage} & \begin{minipage}[b]{\linewidth}\centering
4,5
\end{minipage} \\
\midrule\noalign{}
\endhead
\bottomrule\noalign{}
\endlastfoot
Dice & 0,414 & 0,455 & 0,470 & \textbf{0,489} & 0,484 & 0,426 & 0,332 &
0,272 & 0,210 & 0,167 & 0,127 \\
\end{longtable}
}

Das Optimum liegt mit \(k = 1{,}25\) \textbf{innerhalb} des Rasters (die
Kurve steigt bis 1,25 an und fällt danach ab), ist also kein
Randartefakt der Gittergrenzen.

\textbf{Ergebnis auf dem unabhängigen Testsatz (\(n = 5\),
\(k = 1{,}25\)):}

{\def\LTcaptype{none} % do not increment counter
\begin{longtable}[]{@{}lc@{}}
\toprule\noalign{}
Metrik & Wert \\
\midrule\noalign{}
\endhead
\bottomrule\noalign{}
\endlastfoot
\textbf{Dice} & \textbf{0,513 ± 0,112} \\
IoU & 0,351 ± 0,101 \\
Sensitivität & 0,462 ± 0,215 \\
Spezifität & 0,999 ± 0,001 \\
Precision & 0,736 ± 0,226 \\
\end{longtable}
}

Der voreingestellte Wert \(k = 3{,}0\) (abgeleitet aus der
Normalverteilungs-Annahme \(\mu + 3\sigma \approx 99{,}86\,\%\)) ist für
reale Thermogramme also \textbf{deutlich zu konservativ}. Mit
\(k = 1{,}25\) steigt der Dice-Koeffizient auf einem zuvor ungesehenen
Testsatz von 0,325 auf \textbf{0,513}, wobei die Sensitivität von 0,21
auf 0,46 zunimmt und die Precision erwartungsgemäß von 0,99 auf 0,74
sinkt. Beide Betriebspunkte sind vertretbar: \(k = 3{,}0\) für maximale
Fehlalarmfreiheit, \(k = 1{,}25\) für ein ausgewogenes Verhältnis.

\begin{quote}
\textbf{Einschränkung:} Bei \(n = 5\) Testbildern ist dieses Ergebnis
ein \emph{Hinweis}, kein statistischer Nachweis. Die Standardabweichung
der Sensitivität (± 0,215) zeigt eine erhebliche Streuung zwischen
einzelnen Aufnahmen.
\end{quote}

\begin{center}\rule{0.5\linewidth}{0.5pt}\end{center}

\subsection{5.5 Ablation der geometrischen
Filterregeln}\label{ablation-der-geometrischen-filterregeln}

Die Pipeline verwirft Hotspot-Kandidaten anhand harter geometrischer
Regeln. Zwei davon beruhen auf Annahmen über die Bildkomposition und
könnten echte Befunde auslöschen:

\begin{itemize}
\tightlist
\item
  \textbf{Anatomischer Cutoff:} Komponenten mit Schwerpunkt unterhalb
  von 65 \% der Bildhöhe werden verworfen.
\item
  \textbf{Randfilter:} Komponenten näher als
  \(\max(12\ \text{px};\ 1{,}5\,\%\cdot\min(W,H))\) am Rand werden
  verworfen.
\end{itemize}

Um zu prüfen, ob diese Regeln schaden, wurde ausgezählt, welcher Anteil
der \textbf{annotierten Ground-Truth-Pixel} in die jeweils verworfene
Zone fällt:

{\def\LTcaptype{none} % do not increment counter
\begin{longtable}[]{@{}lc@{}}
\toprule\noalign{}
Filterregel & Anteil der GT-Pixel in der Verwerfungszone \\
\midrule\noalign{}
\endhead
\bottomrule\noalign{}
\endlastfoot
Anatomischer Cutoff (\(y > 0{,}65 \cdot H\)) & \textbf{0,0 \%} \\
Randzone & \textbf{0,0 \%} \\
\end{longtable}
}

Auf diesem Datensatz verwerfen beide Regeln also \textbf{keinen
einzigen} annotierten Befundpixel. Das ist allerdings ausdrücklich
\textbf{keine allgemeine Unbedenklichkeits\-bescheinigung}: Alle 21
Aufnahmen entstanden in derselben Aufnahmegeometrie (Füße mittig im
Bild, gleicher Abstand). Bei abweichender Positionierung -- etwa einem
Fersenulkus im unteren Bilddrittel -- würde der Cutoff den Befund
systematisch unterdrücken. Die Regeln sind damit als
\textbf{datensatzspezifisch} einzustufen und müssten vor einem
Praxiseinsatz entweder entfernt oder an eine automatische
Fußlokalisierung gekoppelt werden.

\begin{center}\rule{0.5\linewidth}{0.5pt}\end{center}

\subsection{5.6 Backend-Paritätstests}\label{backend-parituxe4tstests}

Eine frühere Fassung dieser Arbeit behauptete, Python-, Rust- und
PyTorch-Backend erzeugten „identische Masken''. Diese Aussage war
\textbf{nicht haltbar} und wurde überprüft.

\textbf{Befund der Überprüfung auf allen 21 Realaufnahmen:}

{\def\LTcaptype{none} % do not increment counter
\begin{longtable}[]{@{}lcc@{}}
\toprule\noalign{}
Kennzahl & Vorher & Nach Angleichung \\
\midrule\noalign{}
\endhead
\bottomrule\noalign{}
\endlastfoot
Bitgleiche Masken & 0 / 21 & 0 / 21 \\
Mittlere Masken-IoU (Rust vs.~Python) & 0,29 & \textbf{0,78} \\
Geringste Masken-IoU & 0,11 & 0,48 \\
\end{longtable}
}

Die Ursachenanalyse ergab, dass die Backends \textbf{unterschiedliche
Algorithmen} rechneten: Der Python-Fallback verwendete eine Vorglättung
und eine Multi-Scale-Top-Hat-Transformation über drei Skalen, der
Rust-Core eine einzige Skala ohne Vorglättung. Zusätzlich unterschieden
sich der Otsu-Fallback-Faktor (0,35 vs.~0,30), die
Zusammenhangs\-definition (8er- vs.~4er-Nachbarschaft) und das
morphologische Closing.

Diese Abweichungen wurden im Code beseitigt; die mittlere Masken-IoU
stieg dadurch von 0,29 auf 0,78. \textbf{Exakte Bitgleichheit ist jedoch
prinzipiell nicht erreichbar}, und zwar aus einem grundsätzlichen Grund:
Der Rust-Core nutzt eine \emph{separable} Morphologie nach Lemire mit
rechteckigem Strukturelement (Komplexität \(O(K)\) statt \(O(K^2)\))
sowie eine Chamfer-Approximation der Distanztransformation, OpenCV
dagegen eine andere Approximation und eine abweichende Randbehandlung.
Der Verzicht auf diese Approximationen würde genau den
Geschwindigkeitsvorteil aufheben, für den sie eingeführt wurden.

Die verbleibende Abweichung ist \textbf{systematisch und gerichtet}: Die
Python-Maske ist nahezu eine Teilmenge der Rust-Maske (der Rust-Core
markiert etwas großzügiger). Statt einer unhaltbaren
Gleichheitsbehauptung sichert die Testsuite diese Eigenschaft nun als
überprüfbare Untergrenze ab (\texttt{tests/test\_parity.py}: mittlere
Masken-IoU \(\ge 0{,}70\), Einzelbild \(\ge 0{,}40\)). Die Testsuite
umfasst \textbf{49 Tests}, die vollständig bestehen.

Die \textbf{PyTorch/CUDA-Parität konnte nicht überprüft werden}: Die
verfügbare Grafikkarte (GeForce GTX 1050, Compute Capability 6.1) wird
von der installierten PyTorch-Version (unterstützt ab 7.5) nicht
bedient. Der GPU-Zweig der Testsuite wird daher übersprungen. Aussagen
über GPU-Laufzeiten werden in dieser Arbeit deshalb \textbf{nicht}
getroffen.

\begin{center}\rule{0.5\linewidth}{0.5pt}\end{center}

\subsection{5.7 Ergänzend: synthetische
Entzündungsszenarien}\label{erguxe4nzend-synthetische-entzuxfcndungsszenarien}

Die synthetischen Szenarien aus \texttt{dataset\_evaluator.py} dienen
ausschließlich als Regressionstest der Pipeline, \textbf{nicht} als
Wirksamkeitsnachweis. Der Grund ist methodisch zwingend: Die simulierten
Herde sind gaußförmige Intensitätsmaxima -- also exakt die Signalform,
auf die ein Top-Hat-Filter mit Zirkularitätskriterium ausgelegt ist. Ein
nahezu perfektes Ergebnis ist damit konstruktionsbedingt und nicht
aussagekräftig. Die entsprechenden Kennzahlen finden sich in
\texttt{docs/validation/validation\_report.json}; sie werden hier
bewusst nicht als Gütenachweis in den Vordergrund gestellt.

\begin{center}\rule{0.5\linewidth}{0.5pt}\end{center}

\section{6. Ergebnisdiskussion und Kritische
Würdigung}\label{ergebnisdiskussion-und-kritische-wuxfcrdigung}

\subsection{6.1 Einordnung der Ergebnisse bezüglich der
Arbeitserleichterung}\label{einordnung-der-ergebnisse-bezuxfcglich-der-arbeitserleichterung}

Die Ergebnisse zeigen, dass der Algorithmus schnell rechnet und
auffällige Hitzespitzen mit hoher Treffergenauigkeit lokalisiert: Nahezu
jeder markierte Bildpunkt liegt innerhalb der manuell annotierten Region
(Precision 0,99). Für das Fachpersonal ist genau diese Eigenschaft
entscheidend -- eine Markierung, der man vertrauen kann, verkürzt die
Suchzeit am Bildschirm, während häufige Fehlalarme das Gegenteil
bewirken würden.

Ebenso deutlich ist jedoch zu benennen, was die Ergebnisse
\textbf{nicht} hergeben: Mit einer Sensitivität von 0,21 (bzw. 0,46 bei
angepasstem Schwellenwert) erfasst IGNITE nur einen Teil der auffälligen
Fläche. Eine Entlastung entsteht daher beim \emph{Auffinden} einer
Stelle, nicht bei deren \emph{Beurteilung} oder \emph{Vermessung}. Die
eingangs genannte Zeitersparnis von 3 bis 5 Minuten pro Aufnahme ist
zudem eine eigene Schätzung und wurde nicht in einer Nutzerstudie mit
medizinischem Fachpersonal überprüft -- ein solcher Nachweis wäre für
eine belastbare Aussage zur Arbeitserleichterung erforderlich.

\subsection{6.2 Überprüfung der
Forschungsfragen}\label{uxfcberpruxfcfung-der-forschungsfragen}

\textbf{Forschungsfrage 1 -- teilweise bestätigt.} Der deterministische
Algorithmus schlägt die Otsu-Baseline auf realen Aufnahmen signifikant
(Dice 0,325 vs.~0,004; \(p = 0{,}0046\)). Die ursprünglich formulierte
Zielmarke einer Sensitivität \(> 0{,}95\) wird auf Realaufnahmen jedoch
\textbf{klar verfehlt} (0,206 bei \(k = 3{,}0\); 0,462 bei
\(k = 1{,}25\)). Diese Zielmarke war auf synthetischen Daten formuliert
worden, wo sie trivial erreichbar ist, weil die simulierten Herde exakt
der Filterform entsprechen. Die Selbstkorrektur lautet: \textbf{Auf
synthetischen Daten gemessene Sensitivität ist für dieses Verfahren
keine aussagekräftige Kenngröße.}

\textbf{Forschungsfrage 2 -- widerlegt in der ursprünglichen
Formulierung.} Die Annahme, nativer Rust-Code sei der entscheidende
Geschwindigkeitsfaktor, hat sich nicht bestätigt: OpenCV ist auf
identischem Algorithmus rund 2,4-mal schneller. Die 50-ms-Grenze wird
auf der hochskalierten JPEG-Auflösung verfehlt (86--104 ms), auf der
nativen Sensorauflösung dagegen mit 1,6 ms weit unterschritten. Der
eigentliche Erkenntnisgewinn liegt darin, dass \textbf{die Wahl der
Verarbeitungsauflösung den Laufzeitgewinn dominiert -- nicht die Wahl
der Programmiersprache}. Der Rust-Kern rechtfertigt sich stattdessen
über Abhängigkeitsfreiheit, Determinismus und Speicherbedarf
(\textless{} 25 MB).

\textbf{Forschungsfrage 3 -- bestätigt und in Kapitel 6.3 ausgeführt.}

\subsubsection{Bedeutung des
Robust-MAD-Verfahrens}\label{bedeutung-des-robust-mad-verfahrens}

Das MAD-Verfahren erwies sich bei kühlen Extremitäten als stabiler als
die Schätzung über Mittelwert und Standardabweichung, da einzelne sehr
warme Pixel die Standardabweichung aufblähen und den Schwellenwert
dadurch nach oben verschieben. Zu betonen ist jedoch, dass dieser
Vergleich \textbf{qualitativ} anhand einzelner Aufnahmen erfolgte; eine
quantitative Gegenüberstellung beider Schätzer über den annotierten
Datensatz steht aus.

\begin{center}\rule{0.5\linewidth}{0.5pt}\end{center}

\subsection{6.3 Ausführliche Analyse der Nachteile, Grenzen und
Störfaktoren}\label{ausfuxfchrliche-analyse-der-nachteile-grenzen-und-stuxf6rfaktoren}

\subsubsection{Grenzen der
Datengrundlage}\label{grenzen-der-datengrundlage}

\begin{enumerate}
\def\labelenumi{\arabic{enumi}.}
\tightlist
\item
  \textbf{Sehr kleine Stichprobe.} Ausgewertet wurden \textbf{9
  annotierte Aufnahmen}, davon 5 im unabhängigen Testsatz. Alle
  berichteten Konfidenzintervalle sind entsprechend breit (Dice:
  0,240--0,429). Die Ergebnisse sind als \textbf{Machbarkeitshinweis} zu
  lesen, nicht als statistischer Wirksamkeitsnachweis.
\item
  \textbf{Kein diagnostischer Goldstandard.} Die Referenzmasken wurden
  von \textbf{einer einzigen, medizinisch nicht ausgebildeten Person}
  (dem Autor) unverblindet erstellt. Es liegt weder eine Intra- noch
  eine Inter-Rater-Übereinstimmung vor. Ein Teil der gemessenen
  Abweichung dürfte daher auf die Referenz selbst entfallen, nicht auf
  den Algorithmus.
\item
  \textbf{Keine erkrankten Personen.} Die aufgenommenen Personen sind
  gesund; es wurden thermisch auffällige, nicht pathologisch gesicherte
  Regionen annotiert. Über die Leistung bei echten diabetischen
  Fußläsionen sagt diese Arbeit \textbf{nichts} aus.
\item
  \textbf{Datenverlust während der Arbeit.} Bei der Überprüfung der
  Annotationsdateien stellte sich heraus, dass das Annotationswerkzeug
  beim Weiterblättern bedingungslos speicherte und dadurch sechs zuvor
  erstellte Masken mit leeren Dateien überschrieben hatte. Der Fehler
  wurde behoben (leere Zeichnungen überschreiben keine vorhandene
  Annotation mehr) und die Auswerteroutine prüft Referenzmasken nun auf
  Nicht-Leerheit. Alle vor dieser Korrektur erhobenen
  Ground-Truth-Kennzahlen mussten verworfen werden.
\end{enumerate}

\subsubsection{Physikalische und messtechnische
Grenzen}\label{physikalische-und-messtechnische-grenzen}

\begin{enumerate}
\def\labelenumi{\arabic{enumi}.}
\setcounter{enumi}{4}
\tightlist
\item
  \textbf{Keine Absoluttemperaturen.} Die verwendeten JPEG-Exporte
  enthalten nur relative 8-Bit-Grauwerte einer je Bild dynamisch
  skalierten Palette. Das podiatrische Kriterium \(\Delta T > 2{,}2\,\)K
  (Armstrong et al. 1997) ist damit \textbf{nicht überprüfbar}; die in
  Kapitel 2.2 hergeleitete Radiometrie bleibt theoretisch.
\item
  \textbf{Scheinauflösung des Sensors.} Der Thermosensor liefert
  \(160 \times 120\) Pixel, die Bilddatei suggeriert
  \(1440 \times 1080\). Feine Konturen im Ergebnisbild sind
  Interpolationsprodukte.
\item
  \textbf{Unkontrollierte Aufnahmebedingungen.} Ohne standardisierte
  Akklimatisierung und ohne Protokollierung von Raumtemperatur und
  Luftbewegung sind die Aufnahmen nicht untereinander vergleichbar
  kalibriert.
\item
  \textbf{Störeinflüsse.} Schweiß, Salben, Behaarung sowie die
  Aufnahmeschräge (Lambert-Kosinus-Fehler) verändern die abgestrahlte
  Intensität. JPEG-Blockartefakte können vom Top-Hat-Filter als lokale
  Maxima aufgegriffen werden.
\end{enumerate}

\subsubsection{Algorithmische Grenzen}\label{algorithmische-grenzen}

\begin{enumerate}
\def\labelenumi{\arabic{enumi}.}
\setcounter{enumi}{8}
\tightlist
\item
  \textbf{Systematische Flächenunterschätzung.} Bei einer Precision von
  0,99 und einer Sensitivität von 0,21 markiert IGNITE nur den heißen
  Kern eines Herdes. Für eine Wundflächenvermessung ist das Verfahren
  daher \textbf{ungeeignet}.
\item
  \textbf{Keine Ursachenunterscheidung.} Ein Top-Hat-Filter sucht
  Hitzespitzen. Ob die Wärme von einer Infektion, engen Socken,
  mechanischem Druck oder einer Narbe stammt, kann er prinzipiell nicht
  entscheiden.
\item
  \textbf{Datensatzspezifische Filterregeln.} Der anatomische Cutoff bei
  65 \% der Bildhöhe und der Randfilter verwerfen auf diesem Datensatz
  zwar 0,0 \% der annotierten Pixel (Kapitel 5.5), beruhen aber auf der
  stets gleichen Aufnahmegeometrie. Ein Fersenbefund im unteren
  Bilddrittel würde systematisch unterdrückt -- ein \textbf{blinder
  Fleck}, der bei abweichender Positionierung unbemerkt bliebe.
  \emph{(In Version 5.0 wurde diese Limitation durch vollständige
  Deaktivierung des starren Cutoffs und die Einführung eines gradienten-
  und halobasierten Randschutzes überwunden, wodurch das System
  vollständig universell und anatomie-agnostisch arbeitet.)}
\item
  \textbf{Parameterabhängigkeit.} Der voreingestellte Faktor
  \(k = 3{,}0\) erwies sich als zu konservativ; das auf einem
  Tuning-Satz bestimmte \(k = 1{,}25\) liefert bessere Ergebnisse. Ein
  einzelner globaler Wert kann jedoch nicht jedem Hauttyp und jeder
  Umgebungstemperatur gerecht werden.
\item
  \textbf{Keine Bitgleichheit der Backends.} Python- und Rust-Backend
  liefern trotz Angleichung nur eine mittlere Masken-IoU von 0,78
  (Kapitel 5.6). Anwenderinnen und Anwender erhalten je nach verfügbarem
  Backend leicht abweichende Ergebnisse.
\item
  \textbf{GPU-Pfad unvalidiert.} Der PyTorch/CUDA-Zweig konnte mangels
  kompatibler Hardware nicht getestet werden und ist als
  \textbf{ungeprüft} einzustufen.
\end{enumerate}

\subsubsection{Regulatorische Grenzen}\label{regulatorische-grenzen}

\begin{enumerate}
\def\labelenumi{\arabic{enumi}.}
\setcounter{enumi}{14}
\tightlist
\item
  \textbf{Kein Medizinprodukt.} IGNITE besitzt keine Zertifizierung nach
  EU-MDR und ist ein reiner Forschungsprototyp. Die Software ersetzt
  keine ärztliche Diagnose und darf ausschließlich als
  Orientierungshilfe unter fachlicher Aufsicht verwendet werden.
\end{enumerate}

\begin{center}\rule{0.5\linewidth}{0.5pt}\end{center}

\section{7. Fazit und Ausblick}\label{fazit-und-ausblick}

\subsection{7.1 Gesamtfazit zur
Arbeitswelt-Fragestellung}\label{gesamtfazit-zur-arbeitswelt-fragestellung}

IGNITE zeigt, dass ein vollständig deterministischer, mathematisch
nachvollziehbarer Algorithmus thermisch auffällige Areale auf realen
Aufnahmen signifikant besser markiert als eine einfache
Schwellenwert-Baseline (\(p = 0{,}0046\)), dabei lokal und ohne
Datenübertragung arbeitet und im Millisekundenbereich rechnet.

Der wissenschaftlich ertragreichste Teil dieser Arbeit sind jedoch die
\textbf{widerlegten Ausgangsannahmen}: Der Rust-Kern ist nicht schneller
als OpenCV, die auf synthetischen Daten gemessene Sensitivität von 1,00
hält realen Daten nicht stand, und ein erheblicher Teil der ursprünglich
berichteten Werte beruhte auf methodischen Fehlern (Parameterwahl auf
den Testdaten, überschriebene Referenzmasken, eine nicht existierende
Literaturquelle). Diese Fehler wurden aufgedeckt, korrigiert und
dokumentiert; sämtliche Kennzahlen in Kapitel 5 sind über
\texttt{scripts/run\_validation.py} mit festem Seed reproduzierbar.

Für die Fragestellung der Arbeitswelt bedeutet das: Der Ansatz ist als
\textbf{Orientierungshilfe tragfähig} -- er lenkt den Blick treffsicher
auf die richtige Stelle (Precision 0,99) --, für eine Aussage über die
Ausdehnung eines Befundes oder gar für eine diagnostische Verwendung
reichen die erzielten Werte und die Datengrundlage \textbf{ausdrücklich
nicht aus}.

\subsection{7.2 Zukünftige Erweiterungen für den
Praxiseinsatz}\label{zukuxfcnftige-erweiterungen-fuxfcr-den-praxiseinsatz}

\begin{itemize}
\tightlist
\item
  \textbf{Erweiterung und Absicherung der Referenzdaten:} Annotation der
  verbleibenden 12 Aufnahmen sowie Zweitannotation durch eine
  unabhängige Person, um Cohens \(\kappa\) als Maß der Übereinstimmung
  berichten zu können.
\item
  \textbf{Radiometrische Aufnahmen:} Verwendung des FLIR-Rohdatenformats
  statt der JPEG-Exporte, um mit Absoluttemperaturen zu arbeiten und das
  \(\Delta T > 2{,}2\,\)K-Kriterium (Armstrong et al. 1997) tatsächlich
  prüfen zu können.
\item
  \textbf{Verarbeitung auf nativer Sensorauflösung} (\(160 \times 120\))
  statt auf der hochskalierten JPEG-Ausgabe -- nach den Messungen in
  Kapitel 5.1 der mit Abstand größte Hebel für die Laufzeit.
\item
  \textbf{Ersetzen der starren Geometrieregeln} durch eine automatische
  Fußlokalisierung, um den in Kapitel 6.3 beschriebenen blinden Fleck zu
  beseitigen.
\item
  \textbf{Adaptive Schwellenwertwahl} statt eines globalen Faktors
  \(k\).
\item
  Klinische Validierungsstudie mit fachärztlicher Referenz.
\item
  Automatisch generierter PDF-Befundexport für die Patientenakte.
\end{itemize}

\begin{center}\rule{0.5\linewidth}{0.5pt}\end{center}

\section{8. Literaturverzeichnis}\label{literaturverzeichnis}

(Armstrong et al. 2007) Armstrong, David G. / Holtz-Neiderer, Karin / et
al.: Skin temperature monitoring reduces the risk for diabetic foot
ulceration in high-risk patients. In: \emph{The American Journal of
Medicine}, 2007, Vol. 120, Nr. 12, S. 1042-1046.\\
(Boltzmann 1884) Boltzmann, Ludwig: Ableitung des Stefan'schen Gesetzes,
betreffend die Abhängigkeit der Wärmestrahlung von der Temperatur aus
der electromagnetischen Lichttheorie. In: \emph{Annalen der Physik},
1884, Vol. 258, Nr. 6, S. 291-294.\\
(Stiftung Jugend forscht e. V. 2025) Stiftung Jugend forscht e. V.:
Leitfaden zum Verfassen der schriftlichen Arbeit im Wettbewerb Jugend
forscht. Hamburg, Stand: Juli 2025.\\
(Lemire 2006) Lemire, Daniel: Streaming maximum-minimum filter using no
more than three comparisons per element. In: \emph{Nordic Journal of
Computing}, 2006, Vol. 13, Nr. 4, S. 328-339 (arXiv:cs/0610046).\\
(Jones 1998) Jones, Brian F.: A reappraisal of the use of infrared
thermal image analysis in medicine. In: \emph{IEEE Transactions on
Medical Imaging}, 1998, Vol. 17, Nr. 6, S. 1019-1027. DOI:
10.1109/42.746635.\\
(Steketee 1973) Steketee, J.: Spectral emissivity of skin and
pericardium. In: \emph{Physics in Medicine and Biology}, 1973, Vol. 18,
Nr. 5, S. 686-694. DOI: 10.1088/0031-9155/18/5/307.\\
(Armstrong et al. 1997) Armstrong, David G. / Lavery, Lawrence A. /
Liswood, P. J. / Todd, William F. / Tredwell, Jeffrey A.: Infrared
dermal thermometry for the high-risk diabetic foot. In: \emph{Physical
Therapy}, 1997, Vol. 77, Nr. 2, S. 169-175. DOI: 10.1093/ptj/77.2.169.\\
(Otsu 1979) Otsu, Nobuyuki: A threshold selection method from gray-level
histograms. In: \emph{IEEE Transactions on Systems, Man, and
Cybernetics}, 1979, Vol. 9, Nr. 1, S. 62-66.\\
(Ring and Ammer 2012) Ring, E. Francis J. / Ammer, Kurt: Healthcare
applications of thermal imaging. In: \emph{Physiological Measurement},
2012, Vol. 33, Nr. 3, S. R33-R46.\\
(Stefan 1879) Stefan, Josef: Über die Beziehung zwischen der
Wärmestrahlung und der Temperatur. In: \emph{Sitzungsberichte der
mathematisch-naturwissenschaftlichen Classe der Kaiserlichen Akademie
der Wissenschaften}, 1879, Vol. 79, S. 391-428.

\begin{center}\rule{0.5\linewidth}{0.5pt}\end{center}

\section{9. Unterstützungsleistungen}\label{unterstuxfctzungsleistungen}

Für die Erstellung dieser Arbeit wurden folgende Hilfeleistungen in
Anspruch genommen: * \textbf{Betreuende Lehrkraft / Schule:}
Hilfestellung bei der Formulierung und Durchsicht auf Einhaltung der
Formalien des Jugend forscht Leitfadens. * \textbf{Verwendete
Open-Source-Bibliotheken:} Nutzung freier Programmierwerkzeuge (Rust,
PyO3, Rayon, ndarray, PyTorch, OpenCV, CustomTkinter) gemäß ihren
Open-Source-Lizenzen (MIT / Apache 2.0). * \textbf{Eigenanteil:} Die
Konzeption der Analyse, die Ausarbeitung der Algorithmen, die komplette
Programmierung in Rust und Python, das Interface-Design sowie die
Durchführung aller Tests wurden zu 100 \% eigenständig von mir erbracht.

\protect\phantomsection\label{refs}
\begin{CSLReferences}{1}{1}
\bibitem[\citeproctext]{ref-armstrong2007skin}
Armstrong, David G., Karin Holtz-Neiderer, Christine Wendel, M. Jane
Mohler, Heather R. Kimbriel, and Lawrence A. Lavery. 2007. {``Skin
Temperature Monitoring Reduces the Risk for Diabetic Foot Ulceration in
High-Risk Patients.''} \emph{The American Journal of Medicine} 120 (12):
1042--46.

\bibitem[\citeproctext]{ref-armstrong1997infrared}
Armstrong, David G., Lawrence A. Lavery, P. J. Liswood, William F. Todd,
and Jeffrey A. Tredwell. 1997. {``Infrared Dermal Thermometry for the
High-Risk Diabetic Foot.''} \emph{Physical Therapy} 77 (2): 169--75.
\url{https://doi.org/10.1093/ptj/77.2.169}.

\bibitem[\citeproctext]{ref-boltzmann1884ableitung}
Boltzmann, Ludwig. 1884. {``Ableitung Des Stefan'schen Gesetzes,
Betreffend Die Abh{ä}ngigkeit Der w{ä}rmestrahlung von Der Temperatur
Aus Der Electromagnetischen Lichttheorie.''} \emph{Annalen Der Physik}
258 (6): 291--94.

\bibitem[\citeproctext]{ref-jones1998reappraisal}
Jones, B. F. 1998. {``A Reappraisal of the Use of Infrared Thermal Image
Analysis in Medicine.''} \emph{IEEE Transactions on Medical Imaging} 17
(6): 1019--27. \url{https://doi.org/10.1109/42.746635}.

\bibitem[\citeproctext]{ref-lemire2006streaming}
Lemire, Daniel. 2006. {``Streaming Maximum-Minimum Filter Using No More
Than 3 Comparisons Per Element.''} \emph{Nordic Journal of Computing} 13
(4): 328--39.

\bibitem[\citeproctext]{ref-otsu1979threshold}
Otsu, Nobuyuki. 1979. {``A Threshold Selection Method from Gray-Level
Histograms.''} \emph{IEEE Transactions on Systems, Man, and Cybernetics}
9 (1): 62--66.

\bibitem[\citeproctext]{ref-ring2012healthcare}
Ring, E. Francis J., and Kurt Ammer. 2012. {``Healthcare Applications of
Thermal Imaging.''} \emph{Physiological Measurement} 33 (3): R33--46.

\bibitem[\citeproctext]{ref-stefan1879beziehung}
Stefan, Josef. 1879. {``{Ü}ber Die Beziehung Zwischen Der
w{ä}rmestrahlung Und Der Temperatur.''} \emph{Sitzungsberichte Der
Mathematisch-Naturwissenschaftlichen Classe Der Kaiserlichen Akademie
Der Wissenschaften} 79: 391--428.

\bibitem[\citeproctext]{ref-steketee1973spectral}
Steketee, J. 1973. {``Spectral Emissivity of Skin and Pericardium.''}
\emph{Physics in Medicine and Biology} 18 (5): 686--94.
\url{https://doi.org/10.1088/0031-9155/18/5/307}.

\bibitem[\citeproctext]{ref-jugendforscht2025leitfaden}
Stiftung Jugend forscht e. V. 2025. \emph{Leitfaden Zum Verfassen Der
Schriftlichen Arbeit Im Wettbewerb Jugend Forscht}. Hamburg.

\end{CSLReferences}

\end{document}
